% Volume II: Causal Explosion
% Part III. The Combinatorics of Context
% Chapter 11: Causal Graphs and Expanding Neighborhoods
% Spine: proof-spines/ch011-spine.md

\chapter{Causal Graphs and Expanding Neighborhoods}
\label{ch:ch011}

This chapter begins from the claim that every event is embedded in a causal graph \(G=(V,E)\). To understand a node \(v\), inquiry expands its backward causal neighborhood
\[
B_k(v)=\{u\in V:d(u,v)\leq k\}.
\]
\Definition\ Here \(d(u,v)\) is graph distance and \(k\) the inquiry radius.
If the average branching factor is \(b\), then the size of that neighborhood grows as
\[
|B_k(v)|
\approx
\sum_{j=0}^{k}b^j
=
\frac{b^{k+1}-1}{b-1}.
\]
\ToyModel\ This average-branching approximation makes causal explosion visible before case-specific detail is added.
Even modest values of \(b\) and \(k\) therefore generate an enormous candidate context; this is the elementary mathematics of \textbf{causal explosion}.

The practical question is not whether an investigator can map the whole graph, but how far the relevant neighborhood must expand before attribution is warranted. A visible trigger rarely exhausts the explanation, because background conditions, mediating institutions, incentives, and omissions can all lie inside a larger \(B_k(v)\) and redirect responsibility once they are brought into view. The chapter treats disciplined suspicion as a rule for enlarging adjacency until further expansion stops materially altering the causal account.
